\documentclass[11pt]{article}

\newcommand{\cnum}{Math 245B}
\newcommand{\ced}{Winter 2026}
\newcommand{\ctitle}[4]{\title{\vspace{-0.5in}\cnum, \ced\\Week #1: #2}\author{\vspace{-0.35in}\\#3, #4}}
\usepackage{enumitem}
\usepackage[usenames,dvipsnames,svgnames,table,hyperref]{xcolor}
\usepackage{amsmath, amsfonts}
\usepackage{graphicx} % Required for inserting images
\usepackage{float}
\usepackage{listings}
\usepackage{xcolor}
\usepackage{hyperref}
\usepackage{comment}

\renewcommand*{\theenumi}{\alph{enumi}}
\renewcommand*\labelenumi{(\theenumi)}
\renewcommand*{\theenumii}{\roman{enumii}}
\renewcommand*\labelenumii{\theenumii.}

\definecolor{codegreen}{rgb}{0,0.6,0}
\definecolor{codegray}{rgb}{0.5,0.5,0.5}
\definecolor{codepurple}{rgb}{0.58,0,0.82}
\definecolor{backcolour}{rgb}{0.95,0.95,0.92}
\definecolor{header}{HTML}{205459}
\definecolor{solution}{HTML}{204d52}
\DeclareMathOperator{\spt}{spt}
\DeclareMathOperator{\iter}{int}

\newcommand{\solution}[1]{{{\textcolor{header}{Solution:} \textcolor{solution}{#1}}}}
\setlist[enumerate]{label=(\alph*)}

\lstdefinestyle{mystyle}{
    backgroundcolor=\color{backcolour},
    commentstyle=\color{codegreen},
    keywordstyle=\color{magenta},
    numberstyle=\tiny\color{codegray},
    stringstyle=\color{codepurple},
    basicstyle=\ttfamily\footnotesize,
    breakatwhitespace=false,
    breaklines=true,
    captionpos=b,
    keepspaces=true,
    numbers=left,
    numbersep=5pt,
    showspaces=false,
    showstringspaces=false,
    showtabs=false,
    tabsize=2
}


\begin{document}
\ctitle{\#2}{}{Asher Christian}{006-150-286}
\date{}
\maketitle

\section{Monday}
Assume $H, \langle \cdot, \cdot \rangle$ is a hilbert space.
Definition: Let $I \subset \mathbb{N}$ be a nonempty set and assume that for each $n \in I$,  $E_n$ is a closed vector space.
We say that $H$ is a hilbertian sum of the $(E_n)_{n \in I}$ and we write 
 \[
     H = \bigoplus_{n \in I} E_n
\] 
if the following hold
\begin{enumerate}
    \item If $n,m \in I$ and  $n \ne m$ then  $E_n$ and  $E_m$ are orthogonal
         \[
        \langle u, v \rangle = 0 \quad \forall u \in E_n, v \in E_m
        \] 
    \item The vector space generated by the $E_n$ is dense in  $H$
\end{enumerate}
Theorem: Suppose that $H$ is a hilbertian sum of $(E_n)_{n=1}^{\infty}$ 
\begin{enumerate}
    \item If $u_n \in E_n$ and  $\sum_{n=1}^{\infty}||u_n||^2 < \infty$ then
        \[
            (\sum_{k=0}^{n} u_k)_n
        \] 
        converges to a limit $u \in H$ and  $P_{E_n}(u) = u_n$.
    \item For any  $u \in H$, setting $u+n = P_{E_n}(u)$ we have that
         \[
        \lim_{n\to \infty} ||u - \sum_{k=1}^{n}u_k|| = 0
        \] 
        and we write $u = \sum_{k=1}^{\infty}u_k$
\end{enumerate}
Proof:
\begin{enumerate}
    \item Let $u_n \in E_n$ be such that $\sum_{n=1}^{\infty}||u_n||^2 < +\infty$ set
        \[
        S_n = \sum_{k=1}^{n}u_k
        \] 
        It $n > m \ge 1$ then
         \[
             ||s_n- S_m||^2 = ||\sum_{k=m+1}^{n}u_k||^2 = \sum_{k={m+1}}^{n}||u_k||^2 
        \] 
        since
        \[
            (\sum_{k=0}^{n}||u_k||^2)_n
        \] 
        is a cauchy sequence, we conclude that $(S_n)_n$ is a cauchy sequence.
        Since $H$ is complete $u = \lim_{n\to \infty}S_n$ exists. Fix $n \in \mathbb{N}$ If $m > n$, then 
         \[
             (1) \quad        \langle S_m - u_n, w \rangle, w \in E_n = \langle \sum_{k=1}^{m} u_k - u_n, w \rangle = 0
        \] 
        Letting $m \to +\infty$ we infer $\langle u - u_n, w \rangle = 0$ for all $w \in E_n$ and so  $u_n = P_{E_n}(u)$ since  $E_n$ is a vector space
        \[
        \langle f - P_k(f) , v - P_k(f) \rangle = 0 \quad \forall v \in K
        \] 
        but if $C$ is linear
        \[
        \langle f - P_K(f), v \rangle = 0
        \] 
    \item Let $u \in H$ and set  $u_n = P_{E_n}(u)$ define the operators
        \[
            O_n = \sum_{k=1}^{n}P_{E_k}
        \] 
        note that $O_n(u) = \sum_{k=1}^{n}u_k$.
        Claim
        \[
        ||O_n(u)|| \le ||u||
        \] 
        proof of the claim by characterization of $P_{E_n}(u)$
         \[
        \langle u - u_n, w \rangle = 0 \quad \forall w \in E_n
        \] 
        in particular
        \[
        \langle u - u_n, u_n \rangle = 0
        \] 
        \[
        \langle u , O_n(u) \rangle = \sum_{k=1}^{n}\langle u, u_n \rangle = \sum_{k=1}^{n}||u_n||^2 = || \sum_{k=1}^{n} u_k||^2 = ||O_n(u)||^2
        \] 
        thus
        \[
        ||O_n(u)|| \le ||O_n(u)||||u|| \implies ||O_n(u) || \le ||u||
        \] 
        since the space generated by the $E_n$ is dense in $H$. Given $\epsilon > 0$ we can find $m \in \mathbb{N}$ and $\overline{u} \in span \{E_1,\dots,E_n\}$ such that
        \[
        ||\overline{u} - u|| \le \epsilon
        \] 
        Obeserve that if $n > m$ then  $O_n(u) = \overline{ u}$ %(EXPLAIN THIS)
        and so by claim 1
        \[
        ||O_n(U) - \overline{u}|| = ||O_n(u) - O_n(\overline{ u})|| \le ||\overline{ u} - u || \le \epsilon
        \] 
        thus
         \[
        ||O_n(u) - u|| \le ||O_n(u) - \overline{ u}|| + || \overline{ u} - u|| \le 2\epsilon
        \] 
        for all $\epsilon > 0$ so there exists $m$ such that $n \ge m$ implies  $||O_n(u) - u|| \le 2\epsilon$ and so $\lim_{n\to \infty} O_n(u) =u$
        Remark: Let $u, u_n$ be as in  $(ii)$ of the above theorem then
         \[
        ||u||^2 = ||u - O_n(U) + O_n(u)||^2  = ||u - O_n(u)||^2 + ||O_n(U)||^2 + 2\langle u - O_n(u) , O_n(u) \rangle
        \] 
        \[
         = ||u - O_n(u)||^2 + \sum_{k=1}^{n}||u_k||^2 + 2 \langle u - O_n(u), O_n(u) \rangle
        \] 
        the left and right terms go to zero thus
        \[
        ||u||^2 = \lim_{n\to \infty}\sum_{k=1}^{n} ||u_k||^2
        \] 
\end{enumerate}
Reminder: Riesz representation theorem for hilbert spaces.
For all $L \in H' \exists y \in H$ such that $L(u) = \langle y , u \rangle$ for all $u \in H$.
Furthermore
\[
    ||L||_{H'} = ||y||_{H}
\] 
Avoid to identify $H$ with $H'$\\
Remark: Let $u : (0,1) \to \mathbb{R}$ $C^{1}$ continuous at boundary and set $v = u'$ 
\[
    \int_{0}^{1}u' f = - \int_{0}^{1}uf' \qquad \forall f \in C_0([0,1])
\] 
for all $f$ vanishing at 0 and 1
\[
\int_{0}^{1}vf dx = - \int_{0}^{1}uf'dx \quad \forall f \in C^1_c(0,1)
\] 
Definition: Let $u  \in L^2(0,1)$. Suppose that $v \in L^2(0,1)$ and 
 \[
\int_{0}^{1}vfdx = -\int_{0}^{1}u f'dx \qquad \forall f \in C_C^{1}(0,1)
\] 
we call $v$ the distributional derivative of $u$ and we write
\[
u  \in W^{1,2}(0,1)
\] 
the sobolev space
\[
    W_0^{1,2}(0,1) = \{u \in W^{1,2}(0,1): u(0) = 0, u(1) = 0\}
\] 
the 1 means first derivative the 2 means $L^2$ clearly
\[
W_0^{1,2}(0,1) \subset L^2(0,1)
\] 
is a hilbert space and if $u,f \in W_0^{1,2}(0,1)$ then $\langle u , f \rangle = \langle u, f \rangle_{L^2} + \langle u', f' \rangle_{L^2}$.
Where we use teh distributional derivative.
\[
    (L^2(0,1))' \subset (W_0^{1,2}(0,1))'
\] 
and
\[
W_0^{1,2}(0,1) \subset L^2(0,1)
\] 
if we make the identification $(W_0^{1,2}(0,1))' = W_0^{1,2}(0,1)$ it would mean $L^2(0,1) = W_0^{1,2}(0,1)$
But this is not true so we should not identtify $H$ with $H'$

\section{Wednesday}
Final Exam: Monday March 16, 8:15 AM - 10:45 AM in Boelter 9436
Ex 01: Integration and measure theory, Justify your answer: Has $i$ and $ii$\\
EX 02: $\frac{1}{2}$ about integration and measure theory, other $\frac{1}{2}$ about functional analysis has $i,ii,iii$\\
EX 03: Functional Analysis, has $i,ii$\\
EX 04: Have you filled the evaluation form If not justify your answer, has $i, \to vii$\\

Let $E$ be a banach space and recall that $B_E = \{ x \in E : ||x||_E \le 1\}$.
Notation: If  $f \in E', \alpha \in \mathbb{R}$ and $\epsilon > 0$ we set $V_\epsilon(f,\alpha) = \{ x \in B_E : | f(x) - \alpha| < \epsilon\}$
\[
 = B_E \cap f^{-1}(\alpha - \epsilon, \alpha + \epsilon)
\] 
Theorem: Helly's theorem. Let $f_1,\dots,f_n \in E'$ and $\alpha_1,\dots,\alpha_n \in \mathbb{R}$
then the following are equivalent
\begin{enumerate}
    \item  for every $\epsilon  > 0$ 
        \[
            \bigcap_{i=1}^{n}V(f_i, \alpha_i) \ne \emptyset
        \] 
    \item for all $\beta_1,\dots,\beta_n \in \mathbb{R}$
        \[
            |\sum_{i=1}^{n}\alpha_i \beta_i | \le || \sum_{i=1}^{n} \beta_i f_i ||_{E'}
        \] 
\end{enumerate}
Proof: Assume that $i$ holds and let $\vec \beta = (\beta_1,\dots,\beta_n) \in \mathbb{R}^{n} \setminus \{\vec 0\}$. 
set $S = \sum_{i=1}^{n} |\beta_i| > 0$.
We want to show that for all $\epsilon > 0$
\[
    | \langle \alpha, \beta \rangle_{l_2}| \le \epsilon + ||\sum_{i=1}^{n} \beta f_i ||_E'
\] 
Let $\epsilon > 0$ be arbitrary. Since $i$ holds we can choose $x \in \bigcap_{i=1}^{n}V_\epsilon(f_i, \alpha_i)$.
We have
\[
    |\sum_{i=1}^{n}\alpha_i \beta_i| = |\sum_{i=1}^{n} \beta_i( \alpha_i - f_i(x)) + \sum_{i=1}^{n} \beta_i f_i| \le \sum_{i=1}^{n} |\beta_i| |\alpha_i - f_i(x)| + || \sum_{i=1}^{n} \beta_i f_i ||_{E'}
\] 
since $x \in B_E$ 
\[
\le \epsilon S + ||\sum_{i=1}^{n} \beta_i f_i ||_{E'}
\] 
for all $\epsilon$ which proves $(ii)$.
Assume  $(ii)$ holds. Define $\vec f : E \to \mathbb{R}^{n}$ by
 \[
\vec f(x) = (f_1(x),\dots,f_n(x)) \quad x \in E
\] 
Given $\epsilon > 0$ we want to show that $\bigcap_{i=1}^{n}V_\epsilon(f_i, \alpha_i) \ne \emptyset $.
Set $\vec \alpha = (\alpha_1,\dots,\alpha_n) \in \mathbb{R}^{n}$ we want to show that there exists
$x \in B_E$ such that
\[
    ||\vec f(x) - \alpha||_{L^{\infty}} < \epsilon
\] 
rephrasing, we want to show that  $\vec \alpha \in \overline{ \vec f (B_E)}$.
Assume on the contrary that $\alpha \not\in \overline{ \vec f(B_E)} = K$. Since $K$ is a convex closed set in $ \mathbb{R}^{n}$ 
and $\{ \vec \alpha\}$ is a compact set, by the hahn banach theorem there exists $\vec \beta \ne \vec 0$ and $\gamma \in \mathbb{R}$, $\delta > 0$ such that
\[
\langle \vec \beta, f(x) \rangle + \delta \le \gamma \le \langle \vec \beta, \alpha \rangle - \delta \qquad \forall x \in B_E
\] 
\[
\langle \vec f(x), \beta \rangle \le \gamma - \delta \quad x , -x
\] 
this implies that
\[
|\langle \vec f(x), \beta \rangle | + \delta \le \gamma \le \langle \vec \alpha, \vec \beta \rangle - \delta
\] 
so
\[
| \langle \vec f(x) , \vec \beta \rangle  \le -2 \delta + \langle \vec \alpha, \beta \rangle \qquad \forall x \in B_E
\] 
\[
= |\sum_{}^{}\beta_i f_i(x)|
\] 
so
\[
|| \sum_{i=1}^{n} \beta_i f_i ||_E' \le -2 \delta + \sum_{i=1}^{n} \alpha_i \beta_i \le -2 \delta + | \sum_{i=1}^{n} \alpha_i \beta _i|
\] 
contradicts $(ii)$ so  $\vec \alpha \in \overline{ \vec f(B_E)}$.\\
Theorem: Goldstein Theorem\\
Recall that $J_E: E \to E''$ is defined by  $J_E(f)(x) = f(x) $ for  $x \in E$ and  $f \in E'$
$E$ is reflexive if  $J_E(E) = E''$ which means $J_E(B_E) = B_{E''}$.
Theorem: Goldstein Theorem\\
The set  $J_E(B_E)$ is dense in  $B_{E''}$ for the  $sgim(E'', E')$ topology.
Let  $\xi \in B_{E''}$ we are to show that  $\xi \in \overline{J_E(B_E)}^{\sigma(E'',E')}$ this means we need to show that
if $\xi \in V \in \sigma(E'', E')$ then $V \cap J_E(B_E) \ne \emptyset$ 
Let $f_1,\dots,f_n \in E'$ and $\epsilon > 0$ be such that
\[
    \xi \in W = \bigcap_{i=1}^{n}\{\eta \in E'' : |\eta(f_i) - \xi(f_i)| < \epsilon\}  \qquad \{\eta \in E'': |\xi(f_i) - \eta(f_i)| < \epsilon\}
\] 
set $\alpha_i = \xi(f_i) $ for $i = 1,\dots,n$ showing that $W \cap J_E(B_E) \ne \emptyset$ is equivalent to
showing that there exists  $x \in E$ such that  $|\alpha_i - f_i(x)| < \epsilon$ for all $i$ which means that
 \[
     \bigcap_{i=1}^{n}V_\epsilon(f_i, \alpha_i) \ne \emptyset
\] 
Note that if $\beta_1,\dots,\beta_n \in  \mathbb{R}$ then
\[
    |\sum_{i=1}^{n}\alpha_i \beta_i| = |\sum_{i=1}^{n} \beta_i \xi(f_i)| = | \xi( \sum_{i=1}^{n} \beta_i f_i)| \le || \sum_{i=1}^{n} \beta_i f_i||_{E'}
\] 
since $||\xi||_{E''} \le 1$. By helly's theorem this proves that $\bigcap_{i=1}^{n}V_\epsilon(f_i, \alpha_i) \ne \emptyset$

\section{Friday}
Let $E$ be a Banach space and recall that $S_E = \{x \in E: ||x||_E = 1\}$ and  $B_E == \{x \in E: ||x|| \le 1\}$\\
Milgram-Pettis -  $E$ is uniformly convex implies $E$ is reflexive

Fact: 
\[
J_E : E \to E'' \qquad J_E(f)(x) = f(x) \qquad f \in E', x \in E
\] 
\[
    \overline{J_E(B_E)}^{\sigma(E'',E')} = B_{E''}
\] 
Let $\xi \in S_{E''}$ and  $\delta > 0$
\[
    ||J_E(x) - \xi||_{E''} < \delta
\] 
is false in general
\[
    \sup_{f \in S_{E'}}|J_E(x)(f) - \xi(f)| = ||J_E(x) - \xi||_{E''}
\] 
show that there exists $x \in B_E$ and  $f \in S_{E'}$ such that
\[
    (1) \quad |J_E(x)(f) - \xi(f)| < \frac{\delta}{2} \quad 1 < \xi(f) + \frac{\delta}{2}
\] 
but I can choose $f$ such that 
\[
1 < \xi(f) + \frac{\delta}{2}
\] 
By definition of supremum there exists $f \in S_{E'}$ such that
 \[
     ||\xi||_{E''} < \xi(f) + \frac{\delta}{2}
\] 
use that $f$ and the fact that  $J_E(B_E) \cap \{\eta \in E'': |(\eta - \xi)(f)| < \delta\} \ne \emptyset$

\section{THeorem} Milgram-Pettis\\
If $E$ is uniformly convex then it is reflexive\\
Uniform convex:
For every  $\epsilon > 0$ there exists $\delta \in (0,1)$ such that if $x,y \in B_E$ and  $||x-y|| \ge \epsilon$ then
$|| \frac{x + y}{2} || \le  1 -\delta$.
It suffices to show that $S_{E''} \subset J_E(B_E)$. Let $\xi \in S_{E''}$. Since  $J_E(B_E)$ is closed for the 
strong topology. It suffices to find  $x \in B_E$ such that
$||J_E(x) - \xi||_{E''} < \epsilon$.
Given $\epsilon> 0$ since $E$ is uniformly convex we may choose  $\delta \in (0,1)$ such that
\[
    (2) \quad x,y \in B_E \quad ||x-y||_E \ge \epsilon \implies || \frac{x+y}{2}||_{E} \le 1- \delta
\] 
choose $x \in B_E$ and  $f \in S_{E'}$ such that $(1)$ holds.
Claim:
\[
    ||J_E(x) - \xi||_{E''} \le \epsilon
\] 
Proof: Assume on the contrary that the claim fails so that
$\xi \not\in B_{E''}(J_E(x), \epsilon) = W^{c}$   so $W : = (B_{E''}(J_E(x),\epsilon))^{c}$.
Since $B_{E''}(J_E(x),\epsilon)$ is compact by the banach alaogu-bourbaki theorem $B_{E''}(J_E(x), \epsilon)$ is compact in the $\sigma(E'',E')$ topoogy.
$\xi \in W \in \sigma(E'',E')$ setting $V = \{ \eta \in E'': |(\eta - \xi)(f)| < \frac{\delta}{2} \} \in \sigma(E'', E')$ we have that $\xi \in V \cap W \in \sigma(E'',E')$
By Goldstein Theorem, $V \cap W \cap J_E(B_E) \ne \emptyset$. In other words we can find  $y \in B_E$ such that  $J_E(y) \in V \cap W$.
 $|f(y) - \xi(f) < \frac{\delta}{2}$ and $|f(x) - \xi(f)| < \frac{\delta}{2}$.\\
Hence 
 \[
\delta > |\xi(f) - f(x)| + |\xi(f) - f(y)| \ge \xi(f) - f(x) + \xi(f) - f(y) = 2\xi(f) - f(x+y)
\] 
\[
\xi(f) < \frac{\delta}{2} + f( \frac{x+y}{2})
\] 
 \[
     1< \frac{\delta}{2} + \xi(f) < \delta + f( \frac{x+y}{2})
\] 
by $1$ and we have
\[
    (3) \quad   1 - \delta < f( \frac{x+y}{2}) \le ||\frac{x+y}{2}||_{E}
\] 
But
\[
    (4) \quad    ||x-y||_E = ||J_E(x) - J_E(y)||_{E''} > \epsilon
\] 
since $J_E(y) \in W$ this contradicts the above inequality\\
In conclusion by $(3)$ and  $(4)$
 \[
||x-y||_E < \epsilon \implies ||\frac{x+y}{2}|| > 1 - \delta
\] 
contradicts $(2)$ which proves the claim.
Every Hilbert space is reflexive\\
Proof: Every hilbert space is uniformly convex

\end{document}
